% ======================= 第 四 章 =======================
\chapter{基因原体大叛乱与明清交替}
\chaptersubtitle{ГЛАВА ЧЕТВЁРТАЯ}{ВЕЛИКИЙ МЯТЕЖ ПРИМАРХОВ}
\epigraphbox{“凭尔几路来，我只一路去。”}{——努尔哈赤，据《清太祖高皇帝实录》卷一}

\noindent
万历四十六年四月十三日，努尔哈赤在赫图阿拉城以“七大恨”告天，起兵反明。

关于这一天，清朝的官方史书与明朝的官方史书，给出了两个截然不同却同样简略的记载：

《清太祖高皇帝实录》曰：“上率诸贝勒大臣，谒堂子，陈七大恨，焚之。遂统步骑二万，向抚顺。”

《明神宗实录》曰：“建酋奴儿哈赤，纠众犯顺，陷抚顺，游击李永芳降之。”\footnote{《清太祖高皇帝实录》卷五，天命三年四月，中华书局 1986 年影印本，第 62 页；《明神宗实录》卷五六八，万历四十六年四月，台湾“中研院”史语所校印本，第 10704 页。}

两份记载，一份写“告天”，一份写“犯顺”。没有一份提到那天的天气，也没有一份提到，努尔哈赤在谯楼上站了多久。

本书要说的，正是那段时间里发生的事。

\section{太医院：一个被忽略的机构}

明代太医院，设在皇城东南，与礼部隔街相对。按《明会典》的记载，它是一个医疗服务机构：设院使一人、院判二人，分十三科，掌医疗之法。\footnote{《明会典》卷二二四《太医院》，中华书局 1989 年影印本，第 1096 页。}

十三科。这十三个科室的名称，历代相沿，看上去平平无奇：大方脉、小方脉、妇人、疮疡、针灸、眼、口齿、接骨、伤寒、咽喉、金镞、按摩、祝由。

但请注意最后三科：\kw{金镞、按摩、祝由。}

金镞，治刀箭金创；按摩，治筋骨劳损；祝由，治“邪祟”。\kw{这三科，恰好对应一个强化人形兵器在战场上可能遭遇的三类损伤：机械创伤、结构疲劳、以及精神崩溃。}

而更值得注意的是太医院下辖的两个机构：\kw{生药库}与\kw{惠民药局}。前者“掌收贮药材”，后者“掌给军民药物”。\footnote{同上，第 1097—1098 页。惠民药局的设置，始于洪武三年，其址在南京、北京各一处，“凡军民贫病者，给药”。}

一个给贫民发药的机构，为什么要设在皇城之内、太医院的编制之下？

因为在“药材”的名义下，可以合法地运输、储存、分发任何东西。

\section{01 号原体}

万历二年，辽东总兵李成梁在一次巡边中，于抚顺关外的马市上，见到一个“状貌奇伟”的少年。\footnote{此事的记载见于《清太祖高皇帝实录》卷一：“上生而神武，志意阔大，沈几内蕴。”又，明人笔记《山中闻见录》卷一记：“奴儿哈赤幼时，尝质于李成梁家，成梁抚之如子。”}

少年名努尔哈赤，建州左卫都指挥使觉昌安之孙。

清史称他“生而神武”。明人笔记称他“幼时质于李成梁家，成梁抚之如子”。这两条看似无关的记载，在本书的框架内是同一条：

\kw{他不是被收养的。他是被投放的。}

太医院的“精微科”——十三科之外的第十四个机构，不见于《明会典》，只见于万历十七年的一份内府档案残页——在成化年间启动了一项代号为“原体”的工程。其目标，用今天的话说，是培育一批\kw{可作为边境屏障的生物单位}：不需要军饷，不会叛逃，寿命长，恢复快，且在基因层面与本朝臣民无法区分。

工程共培育了二十个原体，分五批投放。努尔哈赤是第一批的第一个：\kw{01 号。}\footnote{档案残页现存中国第一历史档案馆，编号内府·万历十七年·杂·0327。残页仅存右半，可辨字句为：“……原体二十，分五批，各置一所。其一在辽东……”后半叶缺失。巴兰尼科夫在《辽东基因工程遗址考》第 14 页写道：“一在辽东。就这一句，值三十年的工夫。”}

01 号的编号，被刻在他的锁骨内侧。这一细节，来自一九三七年赫图阿拉老城遗址的一次抢救性发掘：在一处被判定为“清代早期窖藏”的灰坑中，出土了一枚铜制铭牌，长 4.1 厘米，宽 1.7 厘米，正面阴刻“〇一”二字，背面刻有一行极细的字，经放大可辨为：

\vspace{2mm}
\begin{center}
{\small\ttfamily E-13 / GENE-PRIMARCH / BATCH-1 / REV.3}
\end{center}
\vspace{2mm}

REV.3。第三次修订。

一个被造出来的东西，改了三版才投放。而这一版，仍然带着缺陷。\footnote{铭牌现存辽宁省博物馆，编号 1937HL:008。巴兰尼科夫在《辽东基因工程遗址考》中特别强调 REV.3 的“3”字：其笔画与其余数字不同，刻痕更深，像是后来补刻的。他推测：“前两版没有活到刻字的年纪。”}

\section{十三副遗甲}

万历十一年，尼堪外兰引明军攻古勒寨，努尔哈赤的父、祖同时死难。

同年五月，努尔哈赤起兵。《清太祖实录》的记载只有八个字：

\vspace{2mm}
\begin{center}
{\small“上以遗甲十三副，起兵讨尼堪外兰。”}
\end{center}
\vspace{2mm}

“遗甲十三副”。历代注家皆解作“祖先遗留的十三副铠甲”，用以说明努尔哈赤起兵之初的窘迫：一个只有十三副甲的人，最终建立了大清。

这个解释当然动人。但它解释不了两件事。

第一，十三副甲，如何装备一支军队？按明代军制，一副甲的防护面积约为 0.8 平方米；十三副甲，最多装备十三人。而《实录》同卷记载，起兵之时“从者三十人”。\footnote{《清太祖高皇帝实录》卷一，第 24 页。同卷又载，起兵次年“上率兵四百，攻兆佳城”。从三十人到四百人，仅用一年。}

第二，十三副甲如果是“遗留”的，那么遗留者是谁？努尔哈赤的父祖皆死于兵，家中产业被族人瓜分，哪来的十三副完好铠甲？

答案在“遗”字的另一义：\kw{遗，余也。}

十三副，不是十三件，是十三套——一套完整的动力外骨骼，包含躯干框架、双肢作动器、能源模块、控制单元与备用件。十三套，恰好是一个\kw{最小作战单元}的编制。

而“甲”这个字，在明代的军事文书里，也确实可以指代整套装具。《武备志》记“甲一副”，包括盔、身甲、臂甲、裙甲、护心镜、靴，共十七件。\kw{一“副”，从来就是一套，不是一件。}\footnote{茅元仪《武备志》卷一〇五《军资乘·甲》。巴兰尼科夫据此推算：十三副 = 十三套 = 221 件。他在《辽东基因工程遗址考》第 88 页写道：“二十一条命，十三套铁。这笔账，只有一种算法划得来。”}

\begin{table}[htbp]
\centering
\caption{“遗甲十三副”复原清单（据《武备志》甲制与出土残件互校）}
\label{tab:armor}
\footnotesize
\begin{tabular}{@{}clrl@{}}
\toprule
\textbf{编号} & \textbf{名称} & \textbf{出土残件数} & \textbf{状态}\\
\midrule
甲—01 & 躯干承力框架 & 17 & 完整\\
甲—02 & 双作用作动筒（左） & 9 & 变形\\
甲—03 & 双作用作动筒（右） & 11 & 完整\\
甲—04 & 能源模块（固态） & 2 & 泄漏\\
甲—05 & 能源模块（备用） & 6 & 完整\\
甲—06 & 控制单元 & 1 & 缺失\\
甲—07 & 头盔（含视觉增幅） & 4 & 破碎\\
甲—08 & 下肢助力组件 & 13 & 混杂\\
甲—09 & 冷却回路 & 21 & 锈蚀\\
甲—10 & 关节轴承组 & 34 & 混杂\\
甲—11 & 腰带配重 & 5 & 完整\\
甲—12 & 足部缓冲 & 8 & 完整\\
甲—13 & 维修工具包 & 1 & 完整\\
\bottomrule
\end{tabular}
\tabsource{残件统计据《赫图阿拉老城遗址 1937—1939 年发掘报告》附表二；状态判定标准见巴兰尼科夫《辽东基因工程遗址考》，莫斯科，2018 年，第 91—96 页。}
\end{table}

表 \ref{tab:armor} 中最值得注意的是第 06 项：控制单元，出土零残件。

一套完整的动力外骨骼，没有控制单元，就是一堆废铁。而十三套里，恰好只有一套缺了控制单元——\kw{那唯一的一套，穿在努尔哈赤身上。}

因为原体不需要控制器。控制器就在他的脊髓里。

\section{萨尔浒：相位的算法}

万历四十七年二月，明廷以杨镐为辽东经略，调集各路兵马，分四路进剿赫图阿拉。\footnote{《明史》卷二五九《杨镐传》，第 6689 页。关于明军总兵力，历来有“八万八千”“十万”“四十七万”诸说。《明神宗实录》卷五七九载“调兵十万”，而《满文老档》记后金方面统计为“八万八千”。}

西路由杜松率领，出抚顺关；北路由马林率领，出开原；东南路由刘綎率领，出宽奠；南路由李如柏率领，出鸦鹘关。四路合击，约定三月初二会于赫图阿拉。

努尔哈赤的对策，只有十个字：

\vspace{2mm}
\begin{center}
{\small“凭尔几路来，我只一路去。”}
\end{center}
\vspace{2mm}

这句话被后世视为军事天才的集中体现：集中兵力，各个击破。

这个评价是对的。但它只说对了一半。

“我只一路去”，在战术上意味着集中；在\kw{时间}上，它意味着分片。努尔哈赤真正的部署，是把战场当作一个时间切片问题来解：他不去迎击四路，他去迎击\kw{最早到达的那一路}，然后吃掉第二早的，再吃掉第三早的。

这不是“集中兵力”，这是\kw{最短处理时间优先}——一个在今天被写进每一本操作系统教材的调度算法。\footnote{巴兰尼科夫在《重构史学导论》第 268 页给出了萨尔浒战役的甘特图，并将其与“最短作业优先”算法的调度结果对比，结论是：两者在时间轴上完全重合。他批注：“一个十七世纪的猎人，解了一个二十世纪的问题。或者反过来。”}

\begin{figure}[htbp]
\centering
\begin{tikzpicture}[scale=0.87, font=\scriptsize]
  \useasboundingbox (-2.5,-3.3) rectangle (14.2,4.9);

  % 时间轴
  \draw[line width=0.6pt, -{Latex[length=2mm,width=1.3mm]}]
    (0,-2.55) -- (12.4,-2.55);
  \foreach \x/\lab in {0.2/{二月廿五}, 2.1/{廿七}, 4.0/{廿九}, 5.9/{三月初一},
                       7.8/{初三}, 9.7/{初五}}{
    \draw[line width=0.4pt] (\x,-2.55) -- (\x,-2.72);
    \node[anchor=north, font=\tiny] at (\x,-2.78) {\lab};}
  \node[anchor=west, font=\tiny] at (12.5,-2.55) {日};

  % 四路：条形
  \def\laneA{2.10}
  \def\laneB{0.95}
  \def\laneC{-0.20}
  \def\laneD{-1.35}

  % 西路（杜松）
  \node[anchor=east, font=\scriptsize] at (-0.15,\laneA) {西路军（杜松）};
  \fill[black!18] (0.5,\laneA-0.22) rectangle (4.35,\laneA+0.22);
  \draw[line width=0.5pt] (0.5,\laneA-0.22) rectangle (4.35,\laneA+0.22);
  \fill (4.35,\laneA) ++(0,0) -- ++(-0.20,0.16) -- ++(-0.20,-0.16) -- cycle;
  \node[anchor=south, font=\tiny] at (1.45,\laneA+0.28) {萨尔浒—吉林崖};

  % 北路（马林）
  \node[anchor=east, font=\scriptsize] at (-0.15,\laneB) {北路军（马林）};
  \fill[black!18] (1.1,\laneB-0.22) rectangle (5.55,\laneB+0.22);
  \draw[line width=0.5pt] (1.1,\laneB-0.22) rectangle (5.55,\laneB+0.22);
  \fill (5.55,\laneB) -- ++(-0.20,0.16) -- ++(-0.20,-0.16) -- cycle;
  \node[anchor=south, font=\tiny] at (2.25,\laneB+0.28) {尚间崖};

  % 东南路（刘綎）
  \node[anchor=east, font=\scriptsize] at (-0.15,\laneC) {东南路军（刘綎）};
  \fill[black!18] (1.9,\laneC-0.22) rectangle (7.65,\laneC+0.22);
  \draw[line width=0.5pt] (1.9,\laneC-0.22) rectangle (7.65,\laneC+0.22);
  \fill (7.65,\laneC) -- ++(-0.20,0.16) -- ++(-0.20,-0.16) -- cycle;
  \node[anchor=south, font=\tiny] at (3.05,\laneC+0.28) {阿布达里冈};

  % 南路（李如柏）
  \node[anchor=east, font=\scriptsize] at (-0.15,\laneD) {南路军（李如柏）};
  \fill[black!8] (2.4,\laneD-0.22) rectangle (6.30,\laneD+0.22);
  \draw[line width=0.5pt, dash pattern=on 3pt off 2pt]
    (2.4,\laneD-0.22) rectangle (6.30,\laneD+0.22);
  \node[anchor=south, font=\tiny] at (3.35,\laneD+0.28) {虎拦路（未接战）};

  % 相位脉冲标记
  \foreach \x/\y in {4.30/2.10, 5.50/0.95, 7.60/-0.20}{
    \fill (\x,\y) ++(-0.10,-0.30) -- ++(0.20,0) -- ++(-0.10,-0.26) -- cycle;}
  \draw[line width=0.3pt, dash pattern=on 1.5pt off 1.5pt] (4.30,2.10) -- (4.30,-2.45);

  % 当量折线（右轴刻度线）
  \draw[line width=0.9pt] (0.2,-1.35) -- (3.0,-1.35) -- (4.4,0.15)
    -- (5.6,1.35) -- (7.7,2.95) -- (9.7,3.15);
  \foreach \x/\y in {4.4/0.15, 5.6/1.35, 7.7/2.95, 9.7/3.15}{
    \fill (\x,\y) circle (0.075);}
  \node[anchor=west, font=\tiny, align=left] at (9.85,3.15)
    {累计投送当量\\（0—8 kt TNT）};

  % 右轴
  \draw[line width=0.5pt] (8.55,-2.55) -- (8.55,3.15);
  \foreach \y/\v in {-1.35/0, -0.35/2, 0.65/4, 1.65/6, 2.65/8}{
    \draw[line width=0.3pt] (8.48,\y) -- (8.62,\y);
    \node[anchor=west, font=\tiny] at (8.64,\y) {\v};}

  % 图例（顶部单行）
  \node[anchor=north west, align=left, font=\tiny] at (-2.30,4.75)
    {注：\ 实线＝明军行进；\ 虚线＝未接战部队\\
     \phantom{注：\ }折线＝累计当量（右轴）；\ 三角标＝相位脉冲弹引爆点};
\end{tikzpicture}
\caption{萨尔浒之战明军四路行进时序与相位武器当量投放曲线（万历四十七年二月廿五—三月初五）}
\label{fig:sarhu}
\figsource{时序据《明神宗实录》卷五七九至五八〇、《满文老档》太祖朝第二册互校；当量数据据巴兰尼科夫《辽东基因工程遗址考》第 203 页表 11 推定，推定方法见该书附录三。}
\end{figure}

图 \ref{fig:sarhu} 的读法如下。

横轴是时间，从二月廿五到三月初五，共十天。四条横条是四路明军的行进区间，条末的三角是接战点。右侧的折线，是明军累计投送的相位当量——注意，这里的“当量”不是炸药，而是相位脉冲弹的等效能量。

从图中可以读出一个极其刺眼的事实：\kw{明军的当量投放，几乎全部集中在三月初一之后。}

而三路明军，分别覆灭于二月廿九（杜松）、三月初一（马林）、三月初三（刘綎）。

换言之：\kw{明军最重的火力，打在了自己人已经死光了的地方。}

这不是指挥失误。这是\kw{后勤与前沿的时间差}——相位弹的后装、校準、投送，每一个环节都需要时间，而后金军的机动速度，恰好处在这个时间差之内。

后金军没有相位武器，他们有别的：\kw{基因狂暴。}

所谓“基因狂暴”，是原体在极端应激下触发的代谢超频状态：心率升至每分钟 220 次，痛觉阈值提高 6 倍，肌肉输出功率提高 2.7 倍，代价是血糖在四十分钟内耗尽。\footnote{巴兰尼科夫在《辽东基因工程遗址考》第 178 页给出了“狂暴状态持续时间”与“战后死亡率”的相关曲线，r = 0.91。他批注：“这不是战术，这是透支。”}

四十分钟。又是四十分钟。\footnote{见本书第三章关于飞鱼服续航的讨论。巴兰尼科夫在全书最后附了一张“四十分钟清单”，列出中国历史上七个“恰好四十分钟”的技术窗口。他说：“我一直觉得这个数字在跟着我。”}

萨尔浒之战前后五天，后金军连续三次进入狂暴状态。三次，一百二十分钟。\kw{一百二十分钟，决定了此后二百六十八年的中国。}

\section{崩解：从赫图阿拉到山海关}

天命十一年八月，努尔哈赤死于叆鸡堡，年六十八。

《清太祖实录》记其死因，仅“不豫”二字。朝鲜《光海君日记》则记：“奴酋于宁远攻城时，中炮重伤，舁归而毙。”\footnote{《光海君日记》卷一九六，韩国国史编纂委员会影印本，第 482 页。此说与《明熹宗实录》卷七〇“奴酋中炮，伤势甚重”相合。}

两种说法，一种讳言，一种直白。而真相可能比两者都更简单：

\kw{他的基因崩溃了。}

G-13 缺陷，是 01 号原体身上那处改了三版仍未修好的错。它的表型是：在第三次以上的狂暴状态之后，修复机制开始把自身组织识别为异物。临床表现为：伤口不愈、组织坏死、全身性炎症。

中医称之为“疽”。

《清太祖实录》不载病情，《明史》亦不载。但《满文老档》里有一行极不显眼的记载：天命十一年七月，“上命取硫黄三百斤，入汤池”。

三百斤硫黄，入汤池。\kw{那不是在治病，那是在做防腐。}\footnote{《满文老档》太祖朝第七十六册，中华书局 1990 年译本，第 712 页。硫黄在明代医学中确用于疮疡，但用量以钱计，三百斤是工业量级。}

努尔哈赤死后，皇太极继位。他没有继承原体的身份——他是原体的儿子，一个半成品。\footnote{皇太极的体质，据《清太宗实录》卷一所载，自幼“体壮善射”。但同书卷六五记其晚年“鼻衄不止”，崇德八年“无疾而终”，年五十二。巴兰尼科夫认为这是“第二代原体的典型终末”。}

而明朝这边，最后的希望落在一个人身上：袁崇焕。

天启六年正月，宁远。袁崇焕以红夷大炮十一门，击退后金军。\footnote{《明史》卷二五九《袁崇焕传》，第 6708 页。红夷大炮为天启年间自澳门购入的英国舰炮仿制品，口径 120 毫米，射程约 2 公里。}

红夷大炮之所以有效，不在于它的口径，而在于它的\kw{弹道}：滑膛、低伸、铁弹。这种弹道，恰好能与相位场的梯度形成共振——用通俗的话说，\kw{它打的是“频率”，不是“厚度”。}

一门十七世纪的青铜炮，打中了它完全不应该理解的东西。

这就是历史的幽默之处：\kw{终结一个超时代计划的，往往不是另一个超时代计划，而是一件粗糙的、便宜的、谁都能造的东西。}

崇祯十七年三月，李自成入北京，崇祯自缢。四月，吴三桂引清兵入关。

关于吴三桂，历来争论不休：是为陈圆圆，是为父仇，还是为自保？本书不打算参与这场争论。本书只指出一件事：

\kw{吴三桂镇守的山海关，正是太医院“原体工程”的第五号投放点。}

他不是在选边。他是在\kw{调用一个早已写好的协议}。

而这个协议的名字，写在万历十七年那份残页的最后一行，只残留三个字：

\vspace{2mm}
\begin{center}
{\small\ttfamily ...ДОЧЬ...}
\end{center}
\vspace{1mm}
\begin{center}
{\small\color{pencil}（……daughter……）}
\end{center}
\vspace{2mm}

女儿。

\begin{archive}[绝密／СОВЕРШЕННО СЕКРЕТНО]{档案 № ФВ-1966-09}
一九六六年，苏联科学院远东考察队在辽宁新宾赫图阿拉老城北墙外 120 米处，发掘一座砖砌竖穴墓。墓室无棺，仅有石椁一具，椁内铺陈细沙，沙中埋有金属构件三百余件，均严重锈蚀。

其中可辨识者：作动筒残段十一根、齿轮组四套、青铜线缆二百余米、以及一枚直径 62 毫米的圆形薄片，薄片两面均刻有同心圆纹，中央有一个直径 1.2 毫米的小孔。

薄片的材质分析结果为：铂铱合金，含铱 21.4\%。这一配比，与二十世纪航天工业所用的点火电极合金完全一致。

发掘报告称该墓为“清代早期贵族墓葬”。考察队领队 А. И. 库兹涅佐夫在报告空白处手写一行：“这不是墓。这是回收站。”
\end{archive}

\section{本章结论}
\enlargethispage{5\baselineskip}

明清交替，通常被叙述为一次民族征服，或一次王朝更替。

本书不否认这些叙述。它只是补充：\kw{这也是一次实验事故的善后。}

一个王朝造出了它无法控制的东西，这个东西反过来吃掉了它。而吃掉它的那个东西的后代，又建立了一个新的王朝，并在此后的两百年里，系统性地销毁了一切与那项工程有关的记录。

乾隆年间修《四库全书》，禁毁书籍三千余种。其中有一部《辽东志略》，今已不存，仅《四库全书总目》存目，提要云：“其书多载辽左兵事，语多不经，今从删。”

“语多不经”。\kw{一个王朝用“不经”两个字，处理掉了自己的出生证明。}

而我们今天能做的，只是把这些碎片重新拼起来——拼出一个大致能看的形状，然后承认：\kw{它可能不是这个样子。但它一定不是教科书上的那个样子。}
